\begin{figure}[htbp]
\centering
\begin{tikzpicture}[x=1cm,y=1cm]
\node[dissertation block,text width=3.60cm] (base) at (0,3.0)
  {Зафіксована\\конфігурація};

\node[dissertation block,text width=2.85cm] (full) at (-5.1,0)
  {Спільна\\конфігурація};
\node[dissertation block,text width=2.85cm] (minusstr) at (-1.7,0)
  {Без структурних\\ознак};
\node[dissertation block,text width=2.95cm] (minussv) at (1.7,0)
  {Без частотних і\\зорових ознак};
\node[dissertation block,text width=3.05cm] (minusbeh) at (5.1,0)
  {Без правилового\\подієво-поведінкового\\вектора};

\node[dissertation block,text width=12.00cm] (cmp) at (0,-2.8)
  {Агреговане зіставлення шести показників\\за незмінного правила рішення};

% Єдина верхня шина усуває діагональне розгалуження.
\coordinate (configbus) at (0,1.45);
\draw[thick] (base.south) -- (configbus);
\draw[thick] (-5.1,1.45) -- (5.1,1.45);
\draw[dissertation arrow] (-5.1,1.45) -- (full.north);
\draw[dissertation arrow] (-1.7,1.45) -- (minusstr.north);
\draw[dissertation arrow] (1.7,1.45) -- (minussv.north);
\draw[dissertation arrow] (5.1,1.45) -- (minusbeh.north);

% Широкий блок зіставлення приймає чотири окремі вертикальні потоки.
\draw[dissertation arrow] (full.south) -- ([xshift=-5.1cm]cmp.north);
\draw[dissertation arrow] (minusstr.south) -- ([xshift=-1.7cm]cmp.north);
\draw[dissertation arrow] (minussv.south) -- ([xshift=1.7cm]cmp.north);
\draw[dissertation arrow] (minusbeh.south) -- ([xshift=5.1cm]cmp.north);
\end{tikzpicture}
\caption{Схема агрегованого компонентного вилучення}
\label{fig:ablation_design}
\end{figure}
